\section{提案手法}

\subsection{入力と比較範囲}
提案手法は，TPL検出器または関数類似度解析器が出力した関数の組を入力とする静的検証器である．実運用で入力するのは，調査対象のバイナリと，バージョンが既知のTPLから作成した参照バイナリであり，調査対象のソースコードは使用しない．命令列の一致ではなく，ABI，Ghidraが復元した型と格納幅，およびHigh P-codeから得られる振る舞いを比較する．

関数$f$について，2つのバイナリで対応づけた共通の入力範囲を$D_f$とする．各バイナリの引数レジスタやスタック領域を同じ引数へ対応づけ，ABIとGhidraが復元した型・格納幅に基づいて，High P-code上の戻り値格納場所から比較する値を取り出す．比較する振る舞い（以下，観測対象）は，戻り値，仮引数から到達可能なメモリ，大域領域，スレッド局所領域（TLS），ヒープ，\texttt{volatile}アクセス，入出力，ユーザ定義P-code操作，正常・異常終了およびトラップ（trap）とする．\texttt{errno}と浮動小数点環境は既定では対象外であり，並行実行は範囲外である．値，対象領域または効果を比較式で表せない場合は「観測対象外」とせずUNKNOWNにする．入力$x\in D_f$に対する観測結果を$\widehat{O}_i(x)$とすると，本稿では
\[
  \forall x\in D_f,\quad \widehat{O}_1(x)=\widehat{O}_2(x)
\]
を確認できた場合に，2つのバイナリの対象関数を，本稿で定めた観測範囲において等価とする．

異なるABI間では，引数の順序，型，符号およびビット幅を照合する．拡張や切詰めを伴う場合は，変換前後の値の対応と，比較対象としない上位ビットが結果に影響しないことを確認する．入力，戻り値または副作用を同じ対象として対応づけられない場合はUNKNOWNとする．

\subsection{High P-codeからの式抽出と比較}
各関数をGhidraで逆コンパイルし，\texttt{HighFunction}，基本ブロック，\texttt{PcodeOpAST}および\texttt{Varnode}から制御フローとデータ依存を取得する．逆コンパイル結果やP-codeの文字列表現は，演算の意味を決める情報として用いない．RETURNから値の定義を逆向きにたどり，整数演算，比較，分岐条件および\texttt{MULTIEQUAL}を式へ変換する．出力対象のメモリ更新は，2つのバイナリで同じメモリ領域として対応づけた対象と，その書込み範囲を特定できる場合に扱う．

外部関数呼出しは，呼出し先，引数，戻り値および出力に影響するメモリ更新を2つのバイナリの対象関数間で対応づけ，その関係を比較式へ組み込む．メモリは，比較対象とするオブジェクトと書込み順序をHigh P-codeから取得して対応づける．ループは，反復全体の入力と出力の関係を式で表せる場合に比較へ組み込む．CALL，メモリ更新，入出力，終了，トラップ，循環する\texttt{MULTIEQUAL}と反復回数，または演算の関係を確定できない場合はUNKNOWNとする．

2つのバイナリの対象関数から得た終了経路の組合せごとに，経路条件が同時に成立するかを調べる．同時に成立しない組合せを除き，成立し得る組合せについて観測出力が異なる条件をSMTソルバへ与える．EQを出す前には，\tabref{tab:eq-conditions}の条件を確認する．これにより，個々の出力差問い合わせがUNSATでも，CFG上の実行可能性を確認できない経路や，比較式に含まれていない副作用が残る場合をUNKNOWNにできる．解析結果と設定も記録し，異なる実行の結果が混在しないようにする．

\begin{table}[t]
\centering
\caption{EQを判定する前に確認する条件}
\ecaption{Conditions Checked before Returning EQ.}
\label{tab:eq-conditions}
\scriptsize
\setlength{\tabcolsep}{3pt}
\begin{tabular}{@{}p{19mm}p{54mm}@{}}
\hline\hline
条件 & 確認内容 \\
\hline
入力対応 & 引数の意味，型，幅および符号を対応づけられる \\
終了・経路 & 対象とする終了経路を解析できる \\
値・副作用 & 比較対象の戻り値，メモリおよび外部効果を表現できる \\
入力範囲の網羅 & 終了経路の条件が入力領域を覆う \\
出力差 & 到達可能な全経路組で不一致条件がUNSATとなる \\
解析結果の対応 & 比較する結果が，記録された同じ関数対，設定および解析実行に対応する \\
\hline
\end{tabular}
\end{table}

\subsection{判定結果とTPL候補への適用}
判定結果は次の4種類である．定義は本節に集約する．
\begin{description}
  \item[EQ] \tabref{tab:eq-conditions}の全条件を確認し，2つのバイナリで対応づけた同じ入力に対する観測出力の不一致条件がすべてUNSATである．
  \item[NOT\_EQ] 到達可能な入力で観測出力が異なり，その差が比較式に含まれていない値，経路または副作用に依存しない．
  \item[UNKNOWN] EQまたはNOT\_EQの確定に必要な入力対応，経路，値，呼出しまたは副作用を表現または確認できない．
  \item[TIMEOUT] ペア全体の制限時間内に，EQ，NOT\_EQまたはUNKNOWNのいずれかを有効な最終判定として確定できない．
\end{description}
個別のSMT問い合わせが制限時間を超えても，後続の解析によって有効な最終判定を得られた場合はTIMEOUTとしない．対象バイナリの読込み，GhidraまたはSMTソルバの起動，解析中の処理，もしくは判定結果の保存・整合性確認に失敗し，EQ，NOT\_EQ，UNKNOWNまたはTIMEOUTのいずれも有効な最終結果として記録できなかった場合は，ERRORとして別に記録する．NOT\_EQは指定バージョンの候補を除外する根拠となるが，同じTPLの別バージョンまで否定しない．UNKNOWNとTIMEOUTは候補を維持し，後続処理へ理由を渡す．
